\documentclass[12pt]{article}
\usepackage[utf8]{utf8}
\usepackage[spanish]{babel}
\usepackage{geometry}
\geometry{a4paper, margin=2.5cm}
\usepackage{hyperref}
\usepackage{graphicx}
\usepackage{booktabs}
\usepackage{cite}

\title{\textbf{Determinantes Econométricos de la Informalidad Laboral en Ecuador (ENEMDU 2024--2025)}}
\author{\textbf{Matías Carranza}}
\date{\today}

\begin{document}

\maketitle

\begin{abstract}
Este estudio analiza los determinantes socioeconómicos de la informalidad laboral en Ecuador utilizando la base de datos de la ENEMDU. Mediante estimaciones probabilísticas Logit y Probit, se identifican las variables críticas que inciden en la ocupación informal.
\end{abstract}

\vspace{0.5cm}
\noindent\textbf{Repositorio de GitHub:} \url{https://github.com/tu-usuario/tu-repositorio} \\
\noindent\textbf{Dashboard Interactivo (Vercel):} \url{https://tu-dashboard.vercel.app}

\newpage

\section{Introducción}
La informalidad laboral representa un desafío persistente en América Latina \cite{cepal2022}. En Ecuador, las dinámicas del mercado de trabajo muestran patrones específicos que requieren un análisis econométrico riguroso \cite{inec2025}.

\section{Revisión de Literatura}
Estudios recientes señalan que la educación y el nivel salarial son determinantes fundamentales para mitigar la informalidad \cite{oit2023, worldbank2024}. Asimismo, modelos probabilísticos aplicados muestran consistencia en los efectos marginales \cite{gomez2026}.

\section{Metodología}
Se especifican modelos probabilísticos binarios para evaluar la probabilidad de informalidad ($Y_i = 1$):
\begin{equation}
    P(Y_i = 1 | X_i) = F(X_i \beta)
\end{equation}
donde $F(\cdot)$ representa la función de distribución acumulada (logística o normal estándar).

\section{Resultados y Diagnóstico}
Los resultados de los modelos Logit y Probit, junto con sus efectos marginales promedio (AME), coinciden en la dirección de las variables explicativas clave. Las pruebas de capacidad predictiva muestran áreas bajo la curva ROC (AUC) competitivas.

\section{Conclusiones y Recomendaciones}
Se destacan las implicaciones de política pública orientadas a incentivar el empleo formal mediante formación y alivio de cargas administrativas.

\bibliographystyle{apa}
\bibliography{referencias}

\end{document}
